\documentclass[11pt,a4paper]{article}
\usepackage[margin=2.2cm]{geometry}
\usepackage{booktabs, tabularx, array}
\usepackage[hidelinks]{hyperref}
\setlength{\parindent}{0pt}
\setlength{\parskip}{0.5em}
\pagestyle{empty}
% Week 10: why the pooled segment cloglog GLMM suits the who question, in words only.
\renewcommand{\arraystretch}{1.3}

\begin{document}
\sloppy

\subsection*{Why a pooled segment cloglog GLMM}
Geometry and plaque are both measured per segment, but the question is asked per person. The model
therefore learns per segment and answers per person.

\begin{tabularx}{\textwidth}{@{}>{\raggedright\arraybackslash\bfseries}p{3.6cm}X@{}}
\toprule
\textbf{Part of the model} & \textbf{Why it is a good idea here} \\
\midrule
One row per segment & Curvature, taper and plaque are all recorded in each of the nine trunk
  segments. That detail is used instead of being averaged away. \\
cloglog link & Every segment starts plaque-free and is checked once, about ten years later: whether
  plaque appeared is known, but not when. That is a survival question, and the cloglog model is the
  survival model for exactly this kind of data. \\
Added pieces, one test & Segment baseline, SCORE2, segment length and geometry are added together.
  Fitting with and without geometry gives one test of whether geometry adds to SCORE2. \\
Segment baseline & Each of the nine segments gets its
  own starting level of risk: high in the pLAD, where plaque is common, and low in the dLCx, where it is
  rare. The other pieces then only explain differences on top of each segment's normal level, so a
  segment does not look high-risk just because of where it sits. \\
Person effect & Unmeasured risk, such as diet or genes, acts on all segments of a person at once.
  Without it, one person's nine segments would count as nine people and the model would be
  overconfident. \\
Pooling & The question is \emph{who} develops plaque, so the nine segment risks are combined into one
  person risk. With cloglog this is simple: how fast risk builds up in each segment just adds up to how fast it builds up in the person. A
  logistic model has no such simple way to combine them. The person risk comes out at the right level
  and can be compared with SCORE2 alone. \\
\bottomrule
\end{tabularx}

\textbf{An assumption to keep in mind.} Each geometry feature gets one weight for all nine segments,
which keeps the number of weights within what the number of people who developed plaque allows. This
assumes that a bend has the same effect in every segment, which may not be true. A logistic model
would make the same kind of assumption. Whether the effect differs between the LAD, LCx and RCA is
checked by a secondary test.

\end{document}
